% Options for packages loaded elsewhere
\PassOptionsToPackage{unicode}{hyperref}
\PassOptionsToPackage{hyphens}{url}
\documentclass[
]{article}
\usepackage{xcolor}
\usepackage{amsmath,amssymb}
\setcounter{secnumdepth}{-\maxdimen} % remove section numbering
\usepackage{iftex}
\ifPDFTeX
  \usepackage[T1]{fontenc}
  \usepackage[utf8]{inputenc}
  \usepackage{textcomp} % provide euro and other symbols
\else % if luatex or xetex
  \usepackage{unicode-math} % this also loads fontspec
  \defaultfontfeatures{Scale=MatchLowercase}
  \defaultfontfeatures[\rmfamily]{Ligatures=TeX,Scale=1}
\fi
\usepackage{lmodern}
\ifPDFTeX\else
  % xetex/luatex font selection
\fi
% Use upquote if available, for straight quotes in verbatim environments
\IfFileExists{upquote.sty}{\usepackage{upquote}}{}
\IfFileExists{microtype.sty}{% use microtype if available
  \usepackage[]{microtype}
  \UseMicrotypeSet[protrusion]{basicmath} % disable protrusion for tt fonts
}{}
\makeatletter
\@ifundefined{KOMAClassName}{% if non-KOMA class
  \IfFileExists{parskip.sty}{%
    \usepackage{parskip}
  }{% else
    \setlength{\parindent}{0pt}
    \setlength{\parskip}{6pt plus 2pt minus 1pt}}
}{% if KOMA class
  \KOMAoptions{parskip=half}}
\makeatother
\setlength{\emergencystretch}{3em} % prevent overfull lines
\providecommand{\tightlist}{%
  \setlength{\itemsep}{0pt}\setlength{\parskip}{0pt}}
\usepackage{bookmark}
\IfFileExists{xurl.sty}{\usepackage{xurl}}{} % add URL line breaks if available
\urlstyle{same}
\hypersetup{
  pdftitle={Content v2},
  hidelinks,
  pdfcreator={LaTeX via pandoc}}

\title{Content v2}
\author{}
\date{}

\begin{document}
\maketitle

\section{Introduction}\label{introduction}

\begin{itemize}
\tightlist
\item
  Digital Twin definitions → (e.g., Grieves, Tao, Kritzinger\ldots)
\item
  \begin{itemize}
  \tightlist
  \item
    DTs in industrial engineering: why they matter (decision support,
    monitoring, design, troubleshooting).
  \end{itemize}
\item
  Why hybridization is common and beneficial (physics gives structure;
  data gives adaptation).
\item
  The gap: ad hoc hybridization + weak documentation → poor
  transferability.
\item
  \textbf{Contributions (explicit list, 3--5 bullets):}

  \begin{enumerate}
  \tightlist
  \item
    Engineering report of CETIM fluidic-loop DT development
  \item
    DT characterization using Bentley et al.
  \item
    Complementary hybridization perspective (2 dimensions) applied to
    the case
  \item
    Transferable lessons learned + supplementary material
  \end{enumerate}
\end{itemize}

\section{Backgrownd \& terminology}\label{backgrownd-terminology}

Industrial systems are commonly modeled to support prediction,
optimization, and what‑if analysis under operational constraints. Two
broad families of modeling approaches dominate the engineering
literature and practice: \textbf{deductive
(physics‑based/knowledge‑based) models}, which derive system behavior
from first principles, and \textbf{inductive (data‑driven) models},
which infer system behavior from observed data. Increasingly, industrial
digitalization has motivated the \textbf{hybridization} of these
approaches to exploit complementary strengths, interpretability and
extrapolation for deductive models, and flexibility and computational
efficiency for inductive models. In industrial practice, the modeling
pipeline has historically relied on deductive models for simulation and
decision support. However, deductive and inductive approaches provide
complementary capabilities. \textbf{Hybridization} broadly denotes any
systematic strategy that \textbf{combines deductive and inductive
components} within a single modeling workflow to improve qualities such
as accuracy, robustness, computational efficiency, or interpretability.

\subsection{Hybridization at usage (run‑time): coupling deductive and
inductive
models}\label{hybridization-at-usage-runtime-coupling-deductive-and-inductive-models}

When multiple deductive simulators are executed together, each
representing a subsystem, and are numerically coordinated to exchange
variables during execution, the approach is typically described as
\textbf{co‑simulation} (EX1). Co‑simulation thus refers to
\emph{multi‑model} integration where each participating model remains
explicit and mechanistic. A closely related, but distinct, case arises
when \textbf{deductive and inductive models are executed jointly at
run‑time}. In this setting, a physics‑based simulator and a learned
model exchange information during execution, for example:

\begin{itemize}
\tightlist
\item
  {[}Example 1 of hyrbidization at runtime{]}
\item
  {[}Example 2 of hyrbidization at runtime{]}\\
  When the operational output is produced by \textbf{combining
  contributions from both deductive and inductive models during
  deployment}, this constitutes \textbf{hybridization at usage} (EX1).
  The defining feature is that \emph{both} model types are present and
  active in the deployed system, and the final prediction depends on
  their coordinated execution.\\
  To make the concept operational, it is useful to distinguish
  hybridization by \emph{when} and \emph{how} the integration occurs:
  \textbf{hybridization at usage (run‑time integration)} versus
  \textbf{hybridization at creation (design‑time integration)}.
\end{itemize}

\subsection{Hybridization at creation (design‑time): physics‑informed
training and surrogate
modeling}\label{hybridization-at-creation-designtime-physicsinformed-training-and-surrogate-modeling}

Hybridization can also occur during \emph{model construction}, even if
only one model type is used at deployment. A prominent example is
\textbf{surrogate modeling}, where a deductive (physics‑based) model is
used to generate training data for an inductive model that emulates the
simulator's input, output behavior. In this case:

\begin{itemize}
\tightlist
\item
  A deductive model is executed offline across a range of operating
  conditions to produce synthetic data,
\item
  An inductive model is trained on this data to approximate the
  simulator, and
\item
  At deployment, the \textbf{surrogate} (inductive model) may replace
  the original deductive model entirely.\\
  Although the deductive model may no longer be executed at usage time,
  its \textbf{structural assumptions and implied physical relationships
  are transferred} into the learned surrogate through the training
  dataset and learning objective {[}ref{]}. This constitutes
  \textbf{hybridization at creation}, because the integration occurs
  during the development phase rather than during run‑time operation.\\
  \strut \\
\end{itemize}

\section{Fluidic loop case study}\label{fluidic-loop-case-study}

The physical twin considered in this work is a hydraulic loop equipped
with three actuators: (i) pump rotational speed, (ii) opening of a
regulation valve, and (iii) activation of a cooling on/off valve. These
actuators enable the regulation of \textbf{pressure, flow rate, and
temperature} within the loop. The system is instrumented with multiple
sensors (e.g., pressures at key locations, pump speed, and process
temperature) to ensure observability of the operating state. The digital
twin (DT) of this fluidic loop is described following the
characterization framework proposed by Bentley \emph{et al.}, which
structures DTs across \textbf{requirements, realization, deployment, and
operation} phases. This framework is selected because it provides a
formalized and comprehensive description of DT systems, capturing both
their architectural components and operational roles.

\subsection{DT characterization}\label{dt-characterization}

\subsubsection{Requirements, Conceptualization, and
Design}\label{requirements-conceptualization-and-design}

\paragraph{System under study and physical
components}\label{system-under-study-and-physical-components}

The system under study is the fluidic loop itself (MC1), composed of
physical actuators (MC2) and sensing components (MC3). The actuators
include the regulation valve and the cooling unit, while the sensing
layer captures variables such as inlet/outlet pressures, pump speed, and
temperature.

\paragraph{Physical - virtual
interactions}\label{physical---virtual-interactions}

Sensor data are continuously collected and transmitted to the digital
environment through a data infrastructure (data lake), forming the
\textbf{physical-to-virtual interface} (MC4). Conversely, the
\textbf{virtual-to-physical interaction} (MC5) enables the DT to
influence the physical system through actuator commands, such as
adjusting valve opening or managing cooling strategies.

\paragraph{Digital twin services}\label{digital-twin-services}

The DT provides several operational services (MC6), including:

\begin{itemize}
\tightlist
\item
  Cavitation detection,
\item
  \textbf{Virtual sensing}, where flow rate can be estimated in case of
  sensor failure,
\item
  \textbf{Energy consumption analysis}, through simulation of
  alternative cooling strategies,
\item
  \textbf{Divergence detection}, identifying discrepancies between the
  physical and virtual systems to trigger recalibration.
\end{itemize}

\subsubsection{Realization}\label{realization}

The realization of the digital twin (DT) for the fluidic loop relies on
the integration of heterogeneous models, data pipelines, and
computational enablers, orchestrated to support the different
operational services. This section details the DT realization according
to models and data, enabling technologies, system orchestration,
lifecycle evolution, and fidelity considerations.

\paragraph{Digital Twin Models and Data
(MC10)}\label{digital-twin-models-and-data-mc10}

The DT is composed of both \textbf{deductive (physics-based)} and
\textbf{inductive (data-driven)} models, forming a hybrid modeling
architecture.

\begin{itemize}
\item
  \textbf{Physics-based models:}\\
  A 1D thermo-hydraulic simulation model (Simcenter Flomaster) is used
  to represent the fluidic loop. Two variants are employed:

  \begin{itemize}
  \tightlist
  \item
    (i) a simplified hydraulic model, and
  \item
    (ii) a more detailed model incorporating head losses and thermal
    effects.\\
    These models simulate system variables (e.g., pressure, flow rate,
    temperature) based on component parameters and control inputs.
  \end{itemize}
\item
  \textbf{Data-driven models:}\\
  Two inductive models complement the simulation:

  \begin{itemize}
  \tightlist
  \item
    A \textbf{fault localization model} (decision tree), identifying
    which parameters of the physical model are responsible for observed
    discrepancies.
  \item
    A \textbf{fault estimation model} (LS‑SVR), estimating the corrected
    values of these parameters.
  \end{itemize}
\item
  \textbf{Data structure:}\\
  The DT relies on a \textbf{data lake} that stores sensor measurements
  and simulation data. Training datasets for the ML models are generated
  using a combination of:

  \begin{itemize}
  \tightlist
  \item
    sampled control vectors (e.g., valve opening, pump speed),
  \item
    perturbed parameter sets of the simulation model (loss pressur in
    certain pipes),
  \item
    corresponding simulated process variables.
  \end{itemize}
\end{itemize}

\paragraph{Tooling and Enablers (MC11)}\label{tooling-and-enablers-mc11}

The DT implementation is supported by a distributed ecosystem of
industrial and cloud-based tools:

\begin{itemize}
\tightlist
\item
  \textbf{Industrial layer:}

  \begin{itemize}
  \tightlist
  \item
    Beckhoff PLC (automation and control),
  \item
    Industrial PC for local computation.
  \end{itemize}
\item
  \textbf{Interoperability an simulation:}

  \begin{itemize}
  \tightlist
  \item
    FMPy library for executing Functional Mock-up Units (FMU),
  \item
    Model registry for managing simulation assets.
  \end{itemize}
\item
  \textbf{Cloud infrastructure:}

  \begin{itemize}
  \tightlist
  \item
    Azure IoT Hub and Event Hub for data ingestion and streaming,
  \item
    Azure Functions for event-driven computation,
  \item
    Data Lake for storage,
  \item
    OPC UA server for communication.
  \end{itemize}
\item
  \textbf{Application layer:}\\
  Web application interface for visualization and interaction.\\
  These enablers collectively ensure data acquisition, model execution,
  orchestration, and user interaction.
\end{itemize}

\paragraph{Digital Twin Constellation
(MC12)}\label{digital-twin-constellation-mc12}

The DT constellation describes how models, data, and services are
orchestrated within the virtual environment.\\
The architecture follows a \textbf{service-oriented organization}, in
which each DT service (e.g., cavitation detection, virtual sensing,
energy analysis, fault detection) relies on a specific combination of:

\begin{itemize}
\tightlist
\item
  simulation models (FMU-based),
\item
  machine learning components,
\item
  real-time and historical data.\\
  The orchestration is event-driven:
\item
  sensor data streams are ingested into the data infrastructure,
\item
  services are triggered either on-demand or periodically,
\item
  outputs are fed back either to operators (decision support) or to
  actuators.
\end{itemize}

\paragraph{Twinning Process and Digital Twin Evolution
(MC13)}\label{twinning-process-and-digital-twin-evolution-mc13}

The DT evolves over time to improve accuracy and operational
capabilities.

\begin{itemize}
\tightlist
\item
  \textbf{Model evolution:}\\
  The simulation model progresses from a purely hydraulic representation
  to a \textbf{thermo-hydraulic model}, increasing fidelity.
\item
  \textbf{Learning strategy evolution:}\\
  The initial approach relies on generating dedicated datasets and
  training models per use case. This process is expected to evolve
  toward more advanced strategies, such as \textbf{continual learning},
  enabling adaptation without full retraining.
\item
  \textbf{Operational adaptation:}\\
  The DT incorporates mechanisms such as divergence detection, which
  trigger recalibration or parameter updates during operation.
\end{itemize}

\paragraph{Fidelity and Validity Considerations
(MC14)}\label{fidelity-and-validity-considerations-mc14}

The fidelity of the DT depends on the accuracy of both simulation and
data-driven components, as well as their interaction.\\
Key aspects include:

\begin{itemize}
\tightlist
\item
  \textbf{Parameter representation:}\\
  The DT relies on a set of parameters exchanged between the physical
  and virtual systems. Their accuracy, level of abstraction, and
  uncertainty directly influence model validity.
\item
  \textbf{Impact of hybridization:}\\
  Machine learning models (fault localization and estimation) influence
  DT fidelity by dynamically updating simulation parameters. While this
  improves adaptability, it introduces potential uncertainties related
  to model generalization and data quality.
\item
  \textbf{Validation mechanisms:}\\
  Ensuring fidelity requires:

  \begin{itemize}
  \tightlist
  \item
    validation of updated parameters against observed data,
  \item
    monitoring divergence between DT and physical twin,
  \end{itemize}
\end{itemize}

Overall, fidelity is not static but emerges from the \textbf{continuous
interaction between physics-based models, data-driven models, and
real-world observations}. The interaction between these models
illustrates a hybrid architecture, where physics-based simulation is
augmented by data-driven inference mechanisms.\\

\subsubsection{Deployment}\label{deployment}

The deployment of the digital twin (DT) for the fluidic loop is based on
a \textbf{distributed edge--cloud architecture}, ensuring reliable
communication between the physical system, computational components, and
user interfaces. This architecture supports real-time data acquisition,
model execution, and bidirectional interaction with the physical twin.

\paragraph{Digital Twin Technical Connection
(MC15)}\label{digital-twin-technical-connection-mc15}

The DT deployment relies on a multi-layer communication infrastructure
structured along three main segments:

\begin{itemize}
\tightlist
\item
  Factory-to-Edge Layer\\
  At the physical level, sensors and actuators are connected to an
  \textbf{industrial controller (Beckhoff PLC)}, which serves as the
  primary interface with the physical system. This controller aggregates
  sensor measurements (e.g., pressures, temperatures, pump speed) and
  transmits them through an industrial communication stack.\\
  Communication between the factory and the edge is ensured through:

  \begin{itemize}
  \tightlist
  \item
    \textbf{OPC UA protocols}, enabling standardized industrial data
    exchange,
  \item
    \textbf{MQTT-based connectors}, facilitating efficient data
    transmission toward cloud services.
  \end{itemize}
\item
  Edge-to-Cloud Layer\\
  The edge layer communicates with the cloud infrastructure to enable
  scalable data processing and storage. This connection is implemented
  using:

  \begin{itemize}
  \tightlist
  \item
    \textbf{Azure IoT Operations}, acting as a gateway for device
    management and data ingestion,
  \item
    \textbf{Azure Event Hub}, supporting high-throughput streaming of
    telemetry data,
  \item
    Secure communication channels based on \textbf{VPN, firewall
    configurations, and identity management (Azure Active Directory)}.\\
    This layer ensures the reliable transfer of both real-time and
    historical data to the cloud environment, where most DT services are
    executed.
  \end{itemize}
\item
  Cloud-to-Edge Layer\\
  The DT also supports feedback from the virtual environment to the
  physical system. A .NET containerized application, deployed on a
  Kubernetes cluster at the edge (via Azure Arc), interacts with cloud
  APIs to retrieve computed results .\\
  These outputs are transmitted back to the industrial controller using
  the \textbf{TwinCAT ADS communication library}, enabling the Beckhoff
  PLC to actuate system components accordingly.
\end{itemize}

\paragraph{Digital Twin Hosting and Deployment
(MC16)}\label{digital-twin-hosting-and-deployment-mc16}

The DT is hosted across a \textbf{hybrid infrastructure}, combining edge
and cloud resources:

\begin{itemize}
\tightlist
\item
  \textbf{Edge computing layer:}

  \begin{itemize}
  \tightlist
  \item
    Industrial PC for local execution and low-latency operations,
  \item
    Kubernetes cluster for containerized services,
  \item
    Real-time interaction with the physical system.
  \end{itemize}
\item
  \textbf{Cloud computing layer:}

  \begin{itemize}
  \tightlist
  \item
    Azure-based services for scalable data processing,
  \item
    Data Lake for persistent storage,
  \item
    Event-driven computation (Azure Functions),
  \item
    Model management (FMU registry).
  \end{itemize}
\item
  \textbf{Application layer:}\\
  Web-based interface for visualization, monitoring, and user
  interaction.\\
\end{itemize}

\subsubsection{Operation}\label{operation}

The operation phase of the digital twin (DT) describes how the system
delivers value in real-time through continuous interaction with the
physical twin, data processing, and decision support. It focuses on
\textbf{service execution, synchronization mechanisms, and operational
insights} derived from the DT.

\paragraph{Insights and Decision-Making
(MC17)}\label{insights-and-decision-making-mc17}

At the operational level, the DT provides several value-added services
supporting monitoring, diagnosis, and optimization of the fluidic loop:

\subparagraph{\texorpdfstring{\textbf{Cavitation
detection:}}{Cavitation detection:}}\label{cavitation-detection}

The DT evaluates operating conditions to detect cavitation risks within
the hydraulic system using the simulation model.

\subparagraph{Virtual sensing}\label{virtual-sensing}

In case of sensor failure (e.g., flow rate sensor), the DT estimates
missing variables using the simulation model, ensuring continuity of
monitoring.

\subparagraph{Energy consumption
analysis}\label{energy-consumption-analysis}

The DT simulates different cooling strategies (e.g., fast, medium, slow
cooling scenarios) and evaluates their energy impact, enabling informed
operational decisions.

\subparagraph{Divergence detection}\label{divergence-detection}

The DT continuously compares simulated outputs with real measurements to
identify discrepancies. When divergence is detected, it triggers
recalibration or parameter update processes.\\

\subsection{Hybridization perspective}\label{hybridization-perspective}

To ensure consistency with the digital twin (DT) description, the
\textbf{hybridization characteristics (HC)} are introduced by
establishing a direct correspondence with the \textbf{Bentley \emph{et
al.} framework}. While Bentley et al. structure a DT along its lifecycle
dimensions (requirements, realization, deployment, operation), the
hybridization characteristics refine this description by
\textbf{focusing specifically on how deductive and inductive models are
combined}, why this combination is needed, and how it evolves within the
DT. n this perspective, hybridization is not treated as an isolated
concept but as a \textbf{cross-cutting concern} embedded within the DT
characteristic. HC is structured along two complementary dimensions:
\textbf{Motivation} and \textbf{Realization}. This decomposition is
consistent with Bentley \emph{et al.}, where motivation aligns with
\textbf{requirements, conceptualization, and design (MC1--MC6)}, while
realization relates to \textbf{implementation and operation aspects
(MC10--MC14, MC7, MC17)}.

\paragraph{Motivation Dimension}\label{motivation-dimension}

The motivation dimension captures \textbf{why hybridization is
introduced}, independently of its implementation. It formalizes the
problem context, the limitations of existing DT assets, and the expected
capabilities of the hybridized system.

\paragraph{Overview / Context}\label{overview-context}

(Link to MC1: System under study, MC6: DT services, MC10: Models). The
considered DT is based on a \textbf{1D thermo-hydraulic simulation
model} used to support multiple operational services, including flow
simulation, cavitation detection, and energy analysis.\\
This simulation model constitutes the \textbf{core deductive component}
of the DT. However, the increasing demand for adaptive and reliable
services (MC6) highlights the need to extend this model beyond its
original static capabilities. Hybridization is introduced in this
context as a mechanism to \textbf{augment the DT with adaptive
capabilities while preserving the existing simulation-based structure}.

\paragraph{Problem Statement}\label{problem-statement}

(Link to MC6: Services, MC14: Fidelity and validity)\\
The DT faces a fundamental limitation:\\
the \textbf{simulation model is calibrated only once during the creation
phase}, and its parameters remain fixed during operation.\\
As a consequence:

\begin{itemize}
\item
  the DT cannot account for \textbf{system evolution, degradation, or
  unforeseen operating conditions},
\item
  discrepancies arise between simulated outputs and real measurements,
\item
  the \textbf{reliability of DT services (MC6)} is degraded,
\item
  the \textbf{fidelity of the DT (MC14)} decreases over time.\\
  The core problem can therefore be stated as follows:

  \begin{quote}
  \emph{How can the DT maintain alignment with the physical system
  despite evolving conditions, without requiring repeated manual
  recalibration?}
  \end{quote}
\end{itemize}

\paragraph{Existing Twin Assets and
Gaps}\label{existing-twin-assets-and-gaps}

\emph{(Link to MC10: Models and data, MC14: Fidelity)}\\
The DT already includes:

\begin{itemize}
\tightlist
\item
  a \textbf{physics-based simulation model},
\item
  a \textbf{data infrastructure} (sensor streams, data lake),
\item
  an initial \textbf{parameter calibration}.\\
  However, the following gaps are identified:
\item
  \textbf{Static parameterization:} no mechanism for runtime adaptation,
\item
  \textbf{Lack of causal diagnosis:} inability to identify which
  parameters explain observed discrepancies,
\item
  \textbf{Limited robustness:} reduced performance under unmodeled or
  rare conditions.
\end{itemize}

Hybridization addresses these gaps by introducing \textbf{inductive
models} that:

\begin{itemize}
\tightlist
\item
  identify the parameters responsible for deviations,
\item
  infer updated parameter values based on observed data.
\end{itemize}

\paragraph{Intended Deliverable Capability
(IDC)}\label{intended-deliverable-capability-idc}

\emph{(Link to MC6: Services, MC12: DT constellation, MC10: Models)}\\
Hybridization aims to provide the following capabilities within the DT:

\begin{itemize}
\tightlist
\item
  \textbf{Parameter identification capability:}\\
  Determine which model parameters must be updated to explain
  discrepancies.
\item
  \textbf{Parameter estimation capability:}\\
  Infer corrected values or correction laws for the identified
  parameters.
\item
  \textbf{Adaptive simulation capability:}\\
  Enable the simulation model to \textbf{continuously reflect the
  behavior of the physical system} without direct recalibration.\\
  These capabilities enhance the DT by enabling \textbf{adaptive and
  self-consistent behavior}, directly improving service reliability
  (MC6) and model interaction within the DT constellation (MC12).
\end{itemize}

\paragraph{Conceptual Requirements}\label{conceptual-requirements}

\emph{(Link to MC10: Data and models, MC13: Evolution, MC14:
Fidelity)}\\
To ensure the validity of hybridization, the following requirements are
defined:

\begin{itemize}
\tightlist
\item
  \textbf{Coverage of operating conditions:}\\
  Training datasets must span a sufficiently wide range of control
  configurations.
\item
  \textbf{Consistency of dataset semantics:}\\
  Datasets must be generated with fixed nominal parameters; any change
  requires a new dataset version.
\item
  \textbf{Parameter sampling adequacy:}\\
  Each parameter must be sufficiently sampled to allow stable inference.
\item
  \textbf{Validation before deployment:}\\
  Any parameter update must be validated before being applied;
  otherwise, previous values must be retained.\\
  These requirements ensure that hybridization:

  \begin{itemize}
  \tightlist
  \item
    preserves \textbf{model validity and robustness (MC14)},
  \item
    supports controlled \textbf{DT evolution (MC13)}.
  \end{itemize}
\end{itemize}

\paragraph{Deliverable Requirements}\label{deliverable-requirements}

\emph{(Link to MC11: Tooling, MC12: Constellation, MC17: Operation)}\\
The outputs of hybridization namely the inductive models must satisfy
operational constraints:

\begin{itemize}
\tightlist
\item
  \textbf{Low latency:}\\
  Predictions must be computed within time constraints compatible with
  DT services.
\item
  \textbf{Stability:}\\
  Outputs must not disrupt existing simulation workflows.\\
  These requirements ensure seamless integration of hybridization into:
\item
  the DT execution environment (MC11),
\item
  the orchestration of services (MC12),
\item
  real-time operation (MC17).
\end{itemize}

\section{Result}\label{result}

\begin{itemize}
\tightlist
\item
  What becomes clearer when using Bentley + hybridization lens
\item
  Examples of explicit decisions that were previously implicit:

  \begin{itemize}
  \tightlist
  \item
    ??? why ML was used (e.g., compensate unknown friction losses)
  \item
    ??? how it was integrated (residual, surrogate, parameter estimator)
  \end{itemize}
\item
  Evidence types:

  \begin{itemize}
  \tightlist
  \item
    engineering traceability improvements
  \item
    stakeholder feedback (if you have it)
  \item
    reduced ambiguity in documentation

    \begin{itemize}
    \tightlist
    \item
      easier onboarding of new team members (if applicable)
    \end{itemize}
  \end{itemize}
\end{itemize}

\section{Discussion \& Conclusion}\label{discussion-conclusion}

\begin{itemize}
\tightlist
\item
  What your framework enables; next steps.
\end{itemize}

\end{document}
